% Clase 8 — dos espejos perpendiculares: hacen falta TRES imagenes, no una ni dos.
% La cuarta carga (la imagen de las imagenes) es la que restaura la constancia en
% el primer plano despues de agregar la imagen del segundo.
\begin{tikzpicture}[scale=1]

  \fill[figblue!8] (0,0) rectangle (2.6,2.6);

  \draw[figline, line width=1.4pt] (-2.6,0) -- (2.6,0);
  \draw[figline, line width=1.4pt] (0,-2.6) -- (0,2.6);
  \node[figlblaux] at (2.9,-0.22) {$x$};
  \node[figlblaux] at (-0.22,2.85) {$y$};

  % las cuatro cargas
  \node[figpos] at (1.3,1.55) {};
  \node[figlbl, figred] at (1.75,1.78) {$+q$};
  \node[figneg] at (1.3,-1.55) {};
  \node[figlbl, figblue] at (1.78,-1.78) {$-q$};
  \node[figneg] at (-1.3,1.55) {};
  \node[figlbl, figblue] at (-1.78,1.78) {$-q$};
  \node[figpos] at (-1.3,-1.55) {};
  \node[figlbl, figred] at (-1.78,-1.78) {$+q$};

  % las reflexiones, punteadas
  \draw[figaux] (1.3,1.33) -- (1.3,-1.33);
  \draw[figaux] (1.08,1.55) -- (-1.08,1.55);
  \draw[figaux] (-1.3,1.33) -- (-1.3,-1.33);
  \draw[figaux] (-1.08,-1.55) -- (1.08,-1.55);

  \node[figlblaux, align=center, fill=figblue!8, inner sep=1.5pt] at (1.95,0.62)
    {regi\'on donde\\ se resuelve};

  \node[figlbl, align=center] at (0,-3.35)
    {Con dos espejos en \'angulo recto se ven \textbf{cuatro} im\'agenes,
     y eso es exactamente\\[0.2em]
     lo que hace falta: agregar s\'olo una rompe la constancia en el otro plano.};

\end{tikzpicture}
